%%%%%%%%%%%%%%%%%%%%%%%%%%%%%%%%%%%%%%%%%%%%%%%%%%%%%%%%%%%%%%%%%%%%%%%
%                                                                     %
%                        Charte des clubs                             %
%                                                                Beta %
%%%%%%%%%%%%%%%%%%%%%%%%%%%%%%%%%%%%%%%%%%%%%%%%%%%%%%%%%%%%%%%%%%%%%%%
%
%  Auteur :
%    Julien Déoux <julien.deoux@telecomnancy.net>
%    
%  Licence :
%    CC-BY-NC 4.0 (http://creativecommons.org/licenses/by-nc/4.0)
%
%%%



%%%%%%%%%%%%%%%%%%%
%  Configuration  %
%%%%%%%%%%%%%%%%%%%

\documentclass{article} % Jusqu'à preuve du contraire, les documents édités par le
                                                     % Ceten ne font pas 300 pages
\usepackage[a4paper,includeheadfoot,margin=2.54cm]{geometry}
\usepackage[francais]{babel}

\usepackage{hyperref}
\usepackage{graphicx} 
\usepackage{titlesec} 
\usepackage{ifxetex} 
\usepackage[usenames]{xcolor} 
\definecolor{newCeten}{RGB}{130,11,95}
\usepackage{enumitem,amssymb}
\newlist{todolist}{itemize}{2}
\setlist[todolist]{label=$\square$}
\usepackage{multicol}

\ifxetex
	\usepackage{fontspec} 
	\setmainfont{Roboto Slab}
	\setsansfont{Roboto}
	\setmonofont{Roboto Mono}
	\newfontfamily\condensed{Roboto Condensed}
	\newfontfamily\condensedlight{Roboto Condensed Light}
	\newfontfamily\light{Roboto Slab Light}
	\titleformat*{\section}{\Large\condensedlight\color{newCeten}}
\else
	\usepackage[utf8]{inputenc} 
	\usepackage[T1]{fontenc} 
\fi

\title{Charte des clubs}
\author{Julien Déoux}
\date\today

%%%%%%%%%%%%%%
%  Document  %
%%%%%%%%%%%%%%

\begin{document}

	\pagenumbering{roman}

	%---------------%
	% Page de garde %
	%---------------%
	
	\begin{titlepage}
		\begin{center}
			\includegraphics[width=\textwidth]{images/ceten.png}\par
			\vspace{1.5cm}
			{\Huge \light Charte des clubs}\par
			\vspace{1.5cm}
		\end{center}
		\begin{center}
			Club : \underline{\hspace{8cm}}
			Année : \underline{\hspace{1.5cm}}\\
			~\\
			$\square$ Création du club
			\hspace{3cm}
			$\square$ Renouvellement du bureau du club\\
			~\\
			Président du club : \underline{\hspace{7cm}}\\
			Trésorier du club : \underline{\hspace{7cm}}\\
			Secrétaire du club : \underline{\hspace{7cm}}\\
			Référent de TELECOM Nancy : \underline{\hspace{7cm}}\\
		\end{center}
		Type du club :
		\begin{todolist}
			\item Club de services
			\item Club de loisirs
			\item Club évènementiel (Mois de passation :
				\underline{\hspace{5cm}})
		\end{todolist}
		~\\
		Objet du club :\\
		\underline{\hspace{\textwidth}}
		\underline{\hspace{\textwidth}}
		\underline{\hspace{\textwidth}}
		\underline{\hspace{\textwidth}}
		\underline{\hspace{\textwidth}}
 
		\vfill
		\begin{center}
			{\light Ce texte régit le fonctionnement, les obligations et les
			droits attribués aux différents Clubs composant le CETEN, comme
			définis dans les statuts. Un exemplaire est remis au président du
			club, l’autre est archivé au sein du BDE. \\
			Charte modifiée et votée le 12/03/2013 par le BDE du CETEN}
		\end{center}
	\end{titlepage}

	%--------------------%
	% Table des matières %
	%--------------------%
	
	\tableofcontents
	\clearpage

	%----------%
	% Articles %
	%----------%

	\pagenumbering{arabic}

	\begin{multicols}{2}
		
		\section{Définition}
			
			{\small Un "club" est une entité interne au CETEN, composée de
			membres et dédiée à un but ou une activité particulière. Le
			fonctionnement et les engagements pris par le club sont détaillés
			dans cette Charte des Clubs. L’objet spécifique du club est défini
			au début de cette charte.

			Un Club n’a existence que par signature de cette présente charte par
			le Responsable des clubs du BDE, la personne référent de TELECOM
			Nancy et par les président, trésorier et secrétaire du club.

			Faisant partie intégrante de l’association, toutes les règles et
			décisions s’imposant au CETEN s’appliquent également au Club.}

		\section{Type}
			
			{\small Il est défini trois types de Clubs au sein du CETEN :
			\begin{itemize}
				\item Les clubs de type évènementiel, dont le but est
					d’organiser ou de participer à une ou des manifestations
					ponctuelles
				\item Les clubs de type services, qui apportent un service
					particulier dans un domaine précis aux membres du CETEN
				\item Les clubs de type loisirs, regroupant les autres clubs et
					qui sont généralement axés sur le divertissement
			\end{itemize}
			Le type du club doit être défini avec le BDE lors de sa création et
			notifié au début de la charte.}
		
		\section{Durée}
		
			{\small Un club existe jusqu’à la fin de l’année civile spécifiée au
			début de la charte, ou jusqu’à dernier jour du mois de passation
			défini au début de la charte pour les clubs de type évènementiel.
			Cette durée est reconductible par renouvellement du club via cette
			charte.

			Un membre du club l’est jusqu’à la fin du mandat du bureau du club.}
			% À voir

		\section{Composition}
		
			{\small Un club est composé de membres qui doivent être également
			membres du CETEN. Un club se compose au minimum de trois personnes
			distinctes qui sont :
			\begin{itemize}
				\item le président du club, responsable de celui-ci
				\item le trésorier, chargé de tenir régulièrement une
					comptabilité précise du club
				\item le secrétaire, chargé de rédiger les comptes rendus des 
					réunions tenues par le club.
					% Suivi de la liste des membres
			\end{itemize}
			Ces trois personnes déterminent le "bureau restreint" du club, qui a
			pour fonction de gérer l’ensemble du club. Le club n’a lieu
			d’exister sans la présence de ces trois personnes. La structure
			interne du Club est laissée libre au Bureau.
			% Bureau étendu
			Tous les membres du Club devront au préalable avoir fait
			l’inscription à celui-ci via le formulaire fourni en Annexe 1, qui
			sera transmis au responsable Clubs du BDE.
			% À voir
			Chaque bureau de Club doit tenir une liste à jour des membres de son
			club.}
		
		\section{Création}
		
			{\small La création d’un club se fait par trois membres du CETEN
			définissant leur rôle dans ledit club (président, trésorier et
			secrétaire). Le référent de TELECOM Nancy est désigné par la
			Direction des Études. Toute demande de création de club doit être
			oumise au vote du BDE. En cas d’accord et la charte des clubs
			signée, les fondateurs du club deviennent membres du club et
			définissent le bureau du club.}

		Article  6. Activité  au  sein  du  Club   Le  Club  se  limite  aux  activités  décrites  par  l’objet  défini  au   début  de  la  charte  et  uniquement  par  celui-­‐ci.  Il  mettra  en   œuvre  les  moyens  humains  et  logistiques  pour  parvenir  au   mieux  à  la  réalisation  de  ces  activités.  Toute  activité  n’ayant   pas  de  relations  directes  avec  l’objet  du  Club  doit  faire  avant   réalisation  une  demande  auprès  du  BDE.  Le  BDE  aidera,  dans   la   mesure   du   possible   et   du   raisonnable,   le   Club   à   réaliser   ses   activités,   par   un   financement,   du   matériel,   des   formations  et  un  suivi.  
Pour  s’assurer  que  le  Club  exerce  correctement  son  activité,   un  compte-­‐rendu  d’activités  est  demandé  par  le  BDE  tous  les   mois.  
Article  7. Déroulement  des  réunions   En   période   scolaire,   chaque   Club   doit   se   réunir   à   un   minimum  d’une  fois  par  mois  sur  convocation  du  président   ou  à  la  demande  du  tiers  de  ses  membres.  Pour  les  Clubs  de   type   Événementiels,   deux   mois   avant   l’évènement,   cette   période  est  réduite  à  une  fois  tous  les  quinze  jours.  
Seuls  les  membres  du  Club  peuvent  assister  aux  réunions.  Le   bureau  du  Club  se  réserve  le  droit  d’inviter  un  membre  du   CETEN  ou  toute  personne  étrangère  à  l’association  si  elle  est   susceptible  de  l’aider.  Tous  les  membres  du  BDE  sont  invités   permanent  des  réunions  du  Club.  
La   prise   de   décisions   se   fait   à   la   majorité   absolue   des   voix   exprimées.   En   cas   d’égalité,   la   voix   du   président   est   prépondérante.  Toutes  les  délibérations  feront  l’objet  d’une   retranscription   écrite   mise   à   disposition   des   membres   de   l’association.   Si   une   délibération   est   contestée   par   la   majorité  absolue  des  membres  du  Club,  celle-­‐ci  est  annulée   et  doit  être  revue  à  la  prochaine  réunion  du  Club.  
Chaque   réunion   fera   l’objet   d’un   compte-­‐rendu   qui   sera   archivé  par  le  bureau  du  Club  et  transmis  au  BDE  dans  les   trois  jours  suivant  la  réunion.  
Article  8. Propriété  du  Club   Aucun   bien   ne   peut   appartenir   au   Club.   Cependant,   des   biens   sont   confiés   par   le   BDE   au   Club   en   fonction   de   ses  

Chartes  des  Clubs  du  Cercle  des  Étudiants  de  TELECOM  Nancy  du  12/03/2013  
besoins.   Le   président   du   Club   est   responsable   des   biens  
Article  10. Trésorerie   prêtés  par  le  BDE.  La  liste  des  biens  prêtés  devra  être  établie  
Chaque   Club   peut,   s’il   le   souhaite,   posséder   une   caisse   entre   le   Club   et   le   Responsable   Logistique   du   BDE,   et   être  
monnaie.  Le  trésorier  doit  tenir  à  jour  un  journal  de  caisse   mise  à  jour  régulièrement.  
enregistrant   les   entrées/sorties   d’argent   de   la   caisse   du  
Le   BDE   doit,   dans   la   mesure   des   fonds   et   matériels   disponibles,   fournir   au   Club   l’ensemble   des   moyens  
                  Cercle  des  Étudiants  de  TELECOM  Nancy   3/7   Club.  Ce  journal  de  caisse  devra  être  remis  en  fin  d’année  au   trésorier  du  BDE.  
nécessaires   à   la   réalisation   de   ses   activités.   Ces   biens  
Toute   transaction   nécessitant   l’utilisation   du   compte   peuvent   être   utilisés   par   d’autres   membres   du   CETEN   que  
bancaire   du   BDE   (paiement   sur   internet,   émission   ou   ceux  composant  le  Club,  cependant  cette  utilisation  doit  être  
encaissement   d’un   chèque,   etc.)   doit   obligatoirement   être   soumise  à  l’approbation  du  président  du  Club.  
faite  par  le  trésorier  du  BDE.  Néanmoins,  le  trésorier  du  Club  
Le  BDE  se  réserve,  à  titre  exceptionnel,  le  droit  d’utiliser  un   des  biens  confié  au  Club  en  prévenant  son  président.  
indique   dans   sa   comptabilité   la   dépense   ou   la   recette   effectuée  (même  si  celle-­‐ci  provient  ou  est  perçue  par  le  BDE   et  non  par  le  Club).  
Un  Club  peut  également  se  voir  attribuer,  sous  décision  du   BDE,   des   locaux   à   accès   exclusif   ou   partagés   avec   d’autres   Clubs.   Cet   accès   est   réservé   aux   bureaux   des   Clubs   concernés   et   au   BDE.   L’accès   à   ces   locaux   aux   autres   membres  du  Club  est  laissé  à  la  discrétion  des  bureaux  des   Clubs  concernés  et  au  BDE.  
Un  membre  du  Club  peut  fournir  une  avance  pour  un  achat   du  Club.  Le  remboursement  est  validé  et  effectué  par  le  BDE   sous   présentation   d’une   facture   justifiant   la   dépense   occasionnée   et   par   la   signature   de   la   feuille   de   remboursement  fourni  en  annexe.  
Article  9. Comptabilité   Le   président   et   le   trésorier   sont   tenus   responsables   de   la   comptabilité   du   Club.   Ils   doivent   veiller   au   respect   des   budgets  attribués  par  le  BDE  au  début  de  l’année  civile.  
Article  11. Représentation  du  CETEN   Le   Club   n’est   en   aucun   cas   une   entité   morale.   Elle   représente   uniquement   le   CETEN   pour   l’activité   décrite   en   objet   de   la   charte   et   ne   peut   s’accorder   d’aucune   responsabilité  de  l’association.  
Chaque  Club  se  doit  de  préparer  un  budget  prévisionnel  de   son  activité.  Il  doit  se  faire  avant  l’établissement  du  budget   prévisionnel  du  CETEN  de  l’année  à  venir.  Ce  budget  peut-­‐ être   nul   si   le   Club   ne   prévoit   aucune   dépense   et   aucune   recette   durant   l’année.   Il   y   précisera   les   subventions   demandées  au  BDE  et  les  recettes  reversées  à  l’association.  
Toute  demande  de  sponsors,  devis,  accord,  et  autre  contrat   écrit   entre   le   Club   et   une   personne   morale   doit   être   explicitement  faite  au  nom  du  CETEN.  Seules  les  personnes   compétentes   du   BDE   possèdent   les   droits   de   signature.   La   responsabilité  est  alors  portée  sur  l’association.  
Le   budget   prévisionnel   de   chaque   Club   doit   être   nécessairement  équilibré.  
Article  12. Election  du  nouveau  bureau   Tout  Club  est  responsable  du  renouvellement  du  bureau  et  
Une   réunion   avec   l’ensemble   des   Clubs   et   le   BDE   est   organisée  chaque  année  pour  répartir  et  valider  l’ensemble   des  subventions  attribuées  à  chacun  des  Clubs.  
de   son   élection.   Peuvent   se   présenter   au   Bureau   du   Club   uniquement  les  membres  qui  ont  adhérés  au  Club  au  moins   30  jours  avant  l’élection.  Chaque  élection  du  bureau  doit  se   faire  en  présence  d’au  moins  un  représentant  du  BDE.  
A  la  fin  de  son  mandat,  le  Club  devra  présenter  un  budget   final  (y  détaillant  l’ensemble  des  dépenses  et  recettes  du
  
   Club),  pour  pouvoir  l’inclure  à  celui  du  CETEN.  
Le  vote  se  fait  à  main  levée.  Un  vote  à  bulletin  secret  peut   être  admis  si  la  majorité  des  membres  présents  ou  le  Bureau   du  Club  le  souhaite.  Le  vote  par  procuration  est  admis  dans  
L’état  des  dépenses  et  recettes  du  Club  doit  être  géré  par  le  
la  limite  d’une  procuration  par  personne.  
trésorier  du  Club.  Il  tient  informé  régulièrement  le  trésorier   du  BDE  des  flux  d’argent  au  sein  du  Club.  L’utilisation  d’un   échéancier   est   encouragé   pour   savoir   quelles   recettes   et   dépenses  ont  été  réellement  faites.  Chaque  mouvement  doit  
La  passation  doit  se  faire  durant  le  dernier  mois  du  mandat   du  Club  ou  dans  le  mois  de  passation  spécifié  au  début  de  la   charte  pour  les  Clubs  de  type  Évènementiel.  
être  justifié  par  une  pièce  comptable  (factures,  etc.).  Toutes   ces   pièces   comptables   devront   être   archivées   au   sein   du   local  du  Club,  ou  au  BDE  si  le  Club  ne  possède  pas  de  local.   De   plus,   une   copie   numérique   de   ces   pièces   doit   être   transmise  au  BDE.  
Pour   les   Clubs   de   type  Services   et  Évènementiel,   le   scrutin   est  ouvert  à  tous  les  membres  du  CETEN.  Pour  les  Clubs  de   type  Loisirs,  le  scrutin  est  réservé  aux  uniques  membres  du   Club,  qui  ont  dû  adhérer  à  celui-­‐ci  au  moins  30  jours  avant   l’élection.  
L’élection  se  fait  par  poste  (président,  trésorier,  secrétaire)   en  deux  tours.  Pour  chaque  poste,  est  élu  au  premier  tour  la  

Chartes  des  Clubs  du  Cercle  des  Étudiants  de  TELECOM  Nancy  du  12/03/2013  
personne  ayant  obtenu  la  majorité  absolue  des  suffrages.  Si  
membre   de   l’équipe   d’animation   de   TELECOM   Nancy.   Le   aucune  personne  ne  remporte  la  majorité  absolue  des  voix  
référant  a  pour  rôle  de  conseiller  le  bureau  dans  la  gestion   lors   du   premier   tour,   un   second   tour   est   organisé   à   la  
de  son  Club,  dans  ses  activités  et  est  l’interlocuteur  privilégié   majorité  relative  avec  les  deux  personnes  ayant  reçu  le  plus  
pour   les   questions   de   fonctionnement   liées   à   de  voix  au  premier  tour.  En  cas  d’égalité  des  voix,  c’est  au  
l’administration.   Celui-­‐ci   doit   signer   cette   présente   charte   président    sortant  du  Club  de  départager.  
pour   que   l’attribution   des   points   CIPA   puisse   être   validée.   Un  membre  enseignant  ou  administratif  de  TELECOM  Nancy   Le  nouveau  bureau  prend  ses  fonctions,  après  signature  de  
peut  être  référent  d’un  ou  plusieurs  Clubs.  Pour  permettre   la  charte,  le  1er  Janvier  de  l’année  civile  suivante.  Pour  les  
au  référent  de  suivre  son  activité,  le  Club  devra  lui  envoyer   Clubs  de  type  Évènementiel,  la  date  de  prise  de  fonction  est  
un  compte-­‐rendu  d’activité  au  minimum  une  fois  par  mois.   définie   à   la   date   de   signature   de   la   charte   par   le   nouveau   bureau.    
Dans   le   dernier   mois   du   mandat   du   Club,   le   président   du   Club  doit  soumettre  au  référent,  au  directeur  des  études  et   L’élection  du  nouveau  bureau  doit  faire,  comme  pour  toute  
ainsi   qu’au   BDE   un   compte-­‐rendu   de   l’année   du   Club   réunion  du  Club,  l’objet  d’un  Compte-­‐Rendu.  
comprenant  :  
Article  13. Passation  du  Club  
• Un  bilan  moral  ;   Durant   le   mois   suivant   l’élection,   hors   vacances   scolaires,  
• Un  bilan  financier  ;   l’ancien  bureau  du  Club  doit  assister  le  nouveau  bureau  pour  
• La  liste  des  membres  actifs  du  Club  en  y  précisant   sa  prise  de  fonction  future.  
sa  fonction  et  le  travail  réalisé  par  celui-­‐ci  lors  de  
Un  dossier  de  passation  écrit  doit  être  présent  dans  tous  les   Clubs.   Celui-­‐ci   contiendra   toutes   les   informations   nécessaires  pour  la  bonne  continuité  du  Club.  Il  devra  être   enrichi   au   fur   et   à   mesure   par   les   mandats   des   différents   bureaux.  
                  Cercle  des  Étudiants  de  TELECOM  Nancy   4/7   l’année.   De   plus,   le   président   attribue   à   chaque   membre  un  nombre  de  points  CIPA  compris  entre   zéro   et   le   maximum   autorisé   par   le   fascicule   0A   selon   son   implication.   Les   points   du   président   du   Club  sont  attribués  par  le  BDE.  
Article  14. Référent  ALISE  
Le  BDE  peut  demander  une  révision  du  compte-­‐rendu  et  du   nombre   de   points   CIPA   s’il   juge   celui-­‐ci   non   conforme   à   Pour  certains  Clubs  à  grande  importance  au  sein  du  CETEN,  
l’activité  du  Club.   l’association  ALISE  peut  mettre  à  disposition  du  Club  un  ou   plusieurs  référents.  Ayant  déjà  eu  une  expérience  au  sein  de  
Article  16. Perte   de   la   qualité   de   ce   Club,   ils   peuvent   apporter   des   conseils   et   éviter   de  
membre  du  Club   reproduire  les  erreurs  commises  dans  le  passé.  Le  Club  est   tenu   d’échanger   régulièrement   avec   le   référent   ALISE   sur   son  activité.  Tous  les  Comptes  Rendus  des  réunions  devront   également  être  envoyés  au  référent  ALISE.  
La  qualité  de  membre  du  Club  se  perd  par  démission,  perte   de  la  qualité  de  membre  du  CETEN  ou  exclusion  prononcée   par  le  bureau  du  Club  ou  le  BDE.  Une  demande  d’exclusion   d’un  membre  peut  être  demandée  au  BDE  par  les  deux-­‐tiers  
Article  15. Points   CIPA   et   référent   de  
des  membres  du  Club.  
TELECOM  Nancy  
Le  bureau  du  Club  ou  le  BDE  peut  prononcer  l’exclusion  d’un   La  participation  à  un  Club  peut  donner  un  certain  nombre  de  
membre   du   Club   pour   motif   grave.   Le   membre   concerné   points   CIPA   pour   l’obtention   du   diplôme   de   TELECOM  
sera   convoqué   par   ledit   bureau   pour   plaider   sa   cause.   La   Nancy.   Il   en   revient   à   la   direction   de   TELECOM   Nancy   de  
décision  devra  être  approuvée  par  les  deux  tiers  du  bureau.   définir,  dans  le  fascicule  0A,  le  nombre  de  points  maximum  
Le  vote  par  procuration  n’est  pas  permis.   attribués   à   la   participation   du   Club   en   fonction   de   l’importance  du  Club,  de  son  type,  et  de  la  fonction  prise  par  
Si  ce  membre  faisait  partie  du  Bureau,  le  BDE  devra  mettre  
les  membres.  
en   place   au   sein   du   Club   une   élection   pour   désigner   son   remplaçant.   Une   fois   l’élection   faite,   la   Charte   devra   être   Pour   se   voir   attribuer   les   points   CIPA,   le   Club   possède   un  
signée   par   le   nouveau   bureau.   Si   ce   poste   ne   peut   être   référent   pouvant   être   une   personne   enseignante   ou   un  
pourvu   par   une   autre   personne,   le   Club   est   dissout.
     

Chartes  des  Clubs  du  Cercle  des  Étudiants  de  TELECOM  Nancy  du  12/03/2013  
  
Fait  en  deux  exemplaires  à  Villers-­‐lès-­‐Nancy  le  ........................................................................  
Toutes  les  signatures  devront  être  précédées  de  la  mention  «  Lu  et  approuvé  ».  
Le  président  du  Club,   Le  secrétaire  du  Club,  
Le  trésorier  du  Club,   Le  référent  de  TELECOM  Nancy,  
Le  président  du  BDE,   Le  responsable  Clubs  du  BDE,        
                  Cercle  des  Étudiants  de  TELECOM  Nancy   5/7  

Chartes  des  Clubs  du  Cercle  des  Étudiants  de  TELECOM  Nancy  du  12/03/2013  
Annexe  1  :  Bulletin  d’adhésion  au  Club     
CETEN  :  Adhésion  du  Club  ...............................................  pour  l’année  .........................  
Nom  :  ............................................................  Prénom  :  ...........................................................    Année  :  口  1A  口  2A  口  3A  
L’adhérant  au  Club  certifie  être  adhérant  au  CETEN  pour  l’année  courante  et  s’engage  à  respecter  la  Charte  des  Clubs   tenus  à  disposition  par  le  président  du  Club  et  le  responsable  Club  du  BDE.  
Fait  à  .......................................................  le  ..............................            Signature  :  
  
  
Annexe  2  :  Bulletin  de  procuration     
  
CETEN  :  Passation  du  Club  ...............................................  pour  l’année  ...................  
Date  de  la  passation  :  ........./........./.........  
Je  soussigné(e)  ........................................................................................................................................................................,  
  Assisterai  à  la  passation  du  Club  
  N’y   assisterai   pas   et   souhaite   donner   procuration   à   ...........................................................................    
membre  du  CETEN  et  du  Club  pour  me  représenter  lors  du  vote  de  passation  de  ce  Club.  Conformément  à  la  
charte  des  Clubs  limitant  à  un  le  nombre  de  procurations  endossables  par  personne,  je  suis  conscient  que  le  
bureau  du  Club  peut  redistribuer  ma  procuration  en  cas  de  dépassement.  
Fait  à  .......................................................  le  ..............................            Signature  :  
  
     
                  Cercle  des  Étudiants  de  TELECOM  Nancy   6/7  

Chartes  des  Clubs  du  Cercle  des  Étudiants  de  TELECOM  Nancy  du  12/03/2013  
Annexe  3  :  Feuille  de  demande  de  remboursement     
CETEN  :  Feuille  de  demande  de  remboursement  
Je  soussigné(e)  ............................................................................................................................................................,  certifie  avoir  
réglé  la  facture  ci-­‐jointe  du  Club  ......................................................  dont  le  montant  s’élève  à  .....................,...............€  
Fait  à  .......................................................  le  .....................................            Signature  :  
Pièce  justificative  :   Remboursé  par  :  
☐  Ticket  de  Caisse   ☐  Espèces  
☐  Facture  (N°  ..............................)   ☐  Chèque  (N°  ........................)  
☐  Autre  (Motif  :  ....................................................................)   ☐  Virement  Bancaire  (N°  ......................)  
Visa  du  président/trésorier  du  Club,   Visa  du  président/trésorier  du  BDE,  
  
  
  
                  Cercle  des  Étudiants  de  TELECOM  Nancy   7/7  
	\end{multicols}

\end{document}
